\section{Raw-kernel realization theorems}\label{sec:realizations}
\subsection{A controlled binary hidden-state filter}\label{sec:filter}
Let $X_t\in\{0,1\}$ have transition probability $q_0\in(0,1/2)$ of flipping at each raw call. The preparation is $P(X_0=1)=1/2$. A command selects $a$ from a fixed finite subset of $[a_-,a_+]\subset(0,1)$. Conditional on the new hidden state, the report $y\in[-1,1]$ has density
\begin{equation}\label{eq:hmm-kernel}
 g(y\mid X_{t+1}=x,a)=\tfrac12(1+a(2x-1)y).
\end{equation}
Each call costs one acquisition and one unit of physical time. The exploration may be any specified causal command rule, fixed independently of the encoder. A terminal probe reads $X_N$ and ends the experiment. Further filtering reports and this terminal probe are the designated future tests.

\begin{theorem}[Filtering realization]\label{thm:hmm}
For this raw-kernel class, the predictive quotient is the scalar conditional probability $p_t=P(X_t=1\mid H_t)$, with the terminal/continuation interface. It has a common compact invariant interval after initialization. For all $N\ge1$ and $m\ge1$,
\begin{equation}\label{eq:hmm-law}
 R_\cp(N,m)-B_N\asymp R_\on(N,m)-B_N\asymp m^{-2},
 \qquad B_N=\E[p_N(1-p_N)].
\end{equation}
The constants depend only on $q_0,a_-,a_+$, not on $N$ or on the specified command rule. The zero-error upper is an exact Borel finite-label implementation. With rational or certified finite-precision calibration, the separate program in Appendix~\ref{sec:algorithms} has the following nonzero-error guarantee. A uniform update/numerical error $u$ in log-odds distance contributes at most $u/(2q_0)$ to the propagated state error. The constants in this statement are not asserted uniform as $q_0\downarrow0$.
\end{theorem}
\begin{proof}
Direct normalization of \eqref{eq:hmm-kernel} gives, with $r=q_0+(1-2q_0)p$,
\begin{equation}\label{eq:hmm-update}
 f_a(y\mid p)=\tfrac12[1+ay(2r-1)],\qquad
 F_{a,y}(p)=\frac{r(1+ay)}{1+ay(2r-1)}.
\end{equation}
All denominators are at least $1-a_+$. Conditional laws of all future report words and the terminal probe are functions of $p$ by induction. Conversely the terminal probe has mean $p$, so different $p$ are predictively different. This proves quotient identification and minimality without assuming a covariance or information matrix.

Let
\[
 c_- =\frac{q_0(1-a_+)}{1+a_+(1-2q_0)},\qquad c_+=1-c_-.
\]
The interval $D=[c_-,c_+]$ contains $1/2$ and is invariant, because the prediction $r$ always lies in $[q_0,1-q_0]$ and the update is increasing in $r$ and $ay$. Use the metric
$d(p,p')=|\operatorname{logit}(p)-\operatorname{logit}(p')|$.
On $D$ it is equivalent to Euclidean distance, with constants depending only on the displayed parameters. The map $p\mapsto p$ is the task map, and its inverse-coordinate Lipschitz constants follow from the bounded derivative $1/[p(1-p)]$. This provides a one-dimensional global chart and an effective valid domain.

For update stability, put $v=p(1-p)$ and $\rho=1-2q_0$. The derivative of the predicted log-odds with respect to previous log-odds equals
\[
 \frac{\rho v}{q_0(1-q_0)+\rho^2 v}\le\rho.
\]
Indeed $v\le1/4$ and $1-\rho^2=4q_0(1-q_0)$, so the denominator is at least $v$. The observation then adds $\log[(1+ay)/(1-ay)]$, independently of the previous state. Thus every same-report update is $\rho$-Lipschitz in $d$. Proposition~\ref{prop:mechanisms}(a) gives $\Gamma_N\le(2q_0)^{-2}$, including initialization rounding.

It remains to verify actual acquired mass. For each preceding history and command,
\begin{equation}\label{eq:hmm-jacobian}
 \frac{\partial F_{a,y}(p)}{\partial y}
 =\frac{2ar(1-r)}{[1+ay(2r-1)]^2}
 \ge c_y:=\frac{2a_-q_0(1-q_0)}{(1+a_+)^2}>0.
\end{equation}
The conditional density of $y$ is at most $(1+a_+)/2$. Change of variables therefore gives a conditional density of $p_{t+1}$ at most $(1+a_+)/(2c_y)$ on its conditional interval and zero elsewhere. Mixtures over the actual histories and commands preserve this upper bound. Its integral is one. Consequently \eqref{eq:mass} holds on the fixed interval $D$ at every $N\ge1$, with $C_g$ this bound times $|D|$. No uniform posterior law has been inserted. Equations~\eqref{eq:chart} and \eqref{eq:mass} give $\Psi_N(m)=|D|^2/m^2$.

For report continuity, direct integration of the density difference gives
\[
 \norm{f_a(\cdot\mid p)-f_a(\cdot\mid p')}_\TV
 =\frac{a\rho}{2}|p-p'|\le\frac{a_+\rho}{8}d(p,p').
\]
The terminal Bernoulli kernel has total variation $|p-p'|\le d(p,p')/4$. Hence all loss-relevant reports satisfy \textup{(C)}. An equally spaced grid in $p$, the update \eqref{eq:hmm-update}, and exact Borel threshold comparisons define the mathematical label encoder. This is not yet a finite-bit implementation. For rational calibrated constants, Appendix~\ref{sec:algorithms} gives the separate finite-bit routine, including current-report quantization, interval arithmetic and absolute-error conversion. Approximate input and arithmetic with certified absolute errors are absorbed into $u_t$; denominator and log-odds conversion bounds above translate them to the common metric. For rational calibration, fixed-precision interval arithmetic produces any prescribed positive update accuracy; an exact tie need not be resolved, since approximate nearest-grid selection suffices. The persistent datum is the grid index; a bounded-precision register for the current report and arithmetic is temporary, not a retained real state. Summing the geometric error series gives the last error assertion. Apply Theorem~\ref{thm:main} to conclude.
\end{proof}

This is a repeatedly updated nonlinear filter, not a once-prepared finite measurement ensemble. The acquired law is generated by the actual observation density at every positive time. Formula~\eqref{eq:hmm-law} concerns prediction along the selected exploration, not the unrestricted optimal information-gathering controller.

\subsection{Singular geometric acquisition with an absorbing report}\label{sec:sensor}
We next construct a class with no strict contraction during waiting and with genuinely intersecting, different-dimensional acquired supports. Fix finitely many compact boxes $Q_j$ and maps
\[
 z_j:Q_j\longrightarrow[1/4,3/4]^d
\]
that are bi-Lipschitz with uniform constants in physical Euclidean coordinates. They may be nonlinear and their images may intersect. A dimension-zero box is a point. Prepare a latent mark $J_0$ with probabilities $\pi_j$ and, conditionally on $J_0=j$, a point $V$ uniformly on $Q_j$. The hidden target is $Z=z_{J_0}(V)$. After preparation the raw hidden state is $(Z,A_t)$; the latent mark is retained only on the preparation probability space for analyzing its pushforward law. The mark and any zero-width source coordinates are not supplied to the observer.

The apparatus has an unresolved/absorbed flag $A_t\in\{0,1\}$, initially zero. At each raw attempt, when unresolved, an independent coin with success probability $s\in[0,1]$ is tossed. Failure reports $\bot$ and changes nothing. Success reports $Z$ and sets the flag to one. Thereafter every raw attempt reports $\bot$ and preserves $Z$ and the flag. All attempts, including these holds, cost one acquisition and one unit of physical time. The action ``attempt'' is available in both modes, so no unrecorded action-legality bit is assumed. The terminal menu consists of a flag probe returning $A_N$, and $d$ binary probes with success probabilities $Z_i$, each executed only once as a terminal experiment. Give these $d+1$ probes equal weights.

After failure while unresolved, the hidden target distribution is unchanged: failure probability is independent of $Z$. After a successful report the future target probabilities are exactly the reported $Z$, whether or not the latent mark is identifiable. Thus the predictive quotient is the unresolved point $o$, together with the \emph{glued} set of absorbed predictions $z\in\bigcup_jz_j(Q_j)$. At an intersection the mark is quotiented out, not secretly retained. The flag probe distinguishes $o$ from an absorbed state even when their coordinate means coincide.

Set
\begin{equation}\label{eq:sensor-baseline}
 b_* =\frac1{d+1}\sum_{i=1}^d\E[Z_i(1-Z_i)],\quad
 V_* =\frac1{d+1}\sum_{i=1}^d\Var(Z_i),\quad
 h_N=(1-s)^N.
\end{equation}
For $N=0$ take $h_0=1$, including $s=1$. Let $\Psi^{\rm acq}(m)$ be \eqref{eq:psi} for the absorbed component charts with weights $\pi_j$. Widths there are the actual physical chart widths with the fixed task-weight scaling absorbed into the chart constants.

\begin{theorem}[Stratified acquisition law through rank collapse]\label{thm:sensor}
For every $s\in[0,1]$, $N\ge0$, and $m\ge2(J+1)$,
\begin{equation}\label{eq:sensor-law}
 R_\cp(N,m)-b_*\asymp R_\on(N,m)-b_*
 \asymp h_NV_*+(1-h_N)\Psi^{\rm acq}(m).
\end{equation}
The constants depend only on the finite chart and task constants, not on $s,N$, positive widths, or the distance between intersecting strata. When positive widths vanish, zero coordinates are removed from the predictive quotient and the same assertion holds, including point components. If all component score images (including the unresolved point) are separated, Theorem~\ref{thm:main} gives the all-budget version using the full mixture allocation \eqref{eq:psi}, for every $m\ge1$.

The online construction has $\Gamma_N=1$ in the full-mixture transport certificate. It retains no success count or latent component mark. Finite-precision successful reports with metric error at most $\eta$ give a root-mean-square prediction allowance at most $C\sqrt{1-h_N}\eta$, when all hold operations are exact.
\end{theorem}
\begin{proof}
The preparation is the pushforward of the specified box mixture, and the transition is a normalized nonnegative Borel kernel on the compact target space with its two flags. For the success report component one may take the prepared target law as dominating measure; its support is the compact union of the chart images. The success instrument at the unresolved belief has density $s$, and its determining-test numerator has the continuous version $s f(z)$. At absorbed beliefs both densities are zero. The $\bot$ component is atomic. Thus the explicit overwrite below provides exactly the fibre-compatible zero-density extension required in Proposition~\ref{prop:quotient}. The admissible belief domain is the unresolved prepared law together with the point beliefs at absorbed targets. It is a disjoint compact union. All determining future-test evaluations are continuous on it; the explicit overwrite/hold versions below also verify the zero-denominator extension clause of Proposition~\ref{prop:quotient}. The probability of $N$ consecutive failures is exactly $h_N$. Conditional on success by $N$, the target law is still the prepared mixture, because all acquisition coins are independent of it. The acquired predictive law is therefore
\begin{equation}\label{eq:sensor-acquired}
 h_N\delta_o+(1-h_N)\sum_j\pi_j\,(z_j)_\#\operatorname{Unif}(Q_j),
\end{equation}
with the absorbed flag coordinate included in the score. In particular the formal union of all prepared values does not have weight one before it is acquired.

The conditional terminal variance of the flag is zero. On an absorbed history the coordinate variance is $Z_i(1-Z_i)$. On an unresolved history it is
$\E[Z_i(1-Z_i)]+\Var(Z_i)$, by total variance. Averaging proves
\begin{equation}\label{eq:sensor-B}
 B_N=b_*+h_NV_*.
\end{equation}
Conditioned on the raw component mark, the acquired charts have exactly uniform physical box densities. Their bi-Lipschitz constants give \textup{(G)}. Their union can have intersections, different local dimensions, and atomic mass. All these properties follow from the raw preparation and report map.

Equip the quotient with Euclidean distance on its weighted score vector. Report continuity holds as follows. Between two absorbed states the next-attempt kernels are identical, namely $\delta_\bot$. Between unresolved and absorbed modes, the variation distance is at most one, while the flag-score distance is $1/\sqrt{d+1}$; a finite global Lipschitz constant follows. Terminal coordinate kernels are Bernoulli and their variation distance is $|Z_i-Z'_i|$; the flag kernel is bounded by the same mode-separation argument. There is no assertion that distinct Dirac reveal channels are TV-continuous: the reveal action is used only from the single unresolved quotient point. On a successful report define the predictive update to overwrite with that reported absorbed state from every previous representative, including the conditional version at a null report in the absorbed mode. On $\bot$, it is the identity.

The machine initially encodes $o$. At a success it chooses a nearest representative in the union codebook, using the reported physical coordinates; on every failure or later hold it leaves the label unchanged. It never reads $J_0$. Conditional on any raw mark $j$, the success-rounding error is bounded by that chart's cover. Therefore the coefficients are exactly those of Proposition~\ref{prop:mechanisms}(c), with the unresolved point as component zero. This proves $\Gamma_N=1$, irrespective of how many holds occur. It also proves the finite-precision allowance: a single error of size $C\eta$ is introduced on an event of probability $1-h_N$, and none on the complementary event.

For completeness one can see the separated oracle and acquisition terms directly. Lemma~\ref{lem:projection}, conditioned on the success event and its latent mark, and the proof of Corollary~\ref{cor:crossing} give an excess of at least $c(1-h_N)\Psi^{\rm acq}(m)$, even if the decoder also uses centers for the unresolved point. For the upper use one exact unresolved label and $m-1$ absorbed grid labels. Allocation scaling gives
$\Psi^{\rm acq}(m-1)\le C_J\Psi^{\rm acq}(m)$ for $m\ge2(J+1)$: allocating $\lfloor(m-1)/J\rfloor$ to every chart and using the estimate in Corollary~\ref{cor:crossing} proves it. The finite machine just constructed then incurs at most $C(1-h_N)\Psi^{\rm acq}(m)$. Add \eqref{eq:sensor-B}. At $s=0$ or $N=0$ the success term is zero and a single known forecast attains the acquisition term. With zero surviving widths all component charts are points and sufficiently many labels give zero acquired distortion, as the formula states.

For effective implementation, take algebraic calibrated widths and weights, and chart maps with certified finite-precision evaluation. The finite allocation minimum is an algebraic comparison problem. A nearest representative can be selected to within absolute distance $\eta$ by comparing approximated distances; this adds at most $2\eta$ to the best covering radius. No exact comparison of arbitrary input reals is required. A finite readonly table of approximated representatives or a grid-generation program, with separately charged description and temporary workspace, suffices. Thresholds are compared in physical coordinates; no uniform error estimate divides by a collapsing width. If the unresolved target mean also requires numerical integration, its certified score error $\eta_0$ is charged separately; the total readout allowance is then at most $\sqrt{h_N}\eta_0+C\sqrt{1-h_N}\eta$. A table or algorithm for this mean is part of the supplied effective calibration, not a free integration oracle. Appendix~\ref{sec:algorithms} distinguishes these costs. This proves the claimed implementation statements.
\end{proof}

\paragraph{A concrete singular family.}
In two coordinate dimensions take equal mixture weights on
\[
 z_1(u)=(1/2+u,1/2),\quad |u|\le a/2,
 \qquad
 z_2(v,w)=(1/2+v,1/2+w),\quad |v|\le b/2,\ |w|\le c/2,
\]
where $0\le a,b,c\le1/4$. The actual acquired law has a one-dimensional singular component on a line and a two-dimensional component on a rectangle. The former lies inside the latter when their horizontal ranges overlap. If $c=0$, the rectangle becomes a line; if also $b=0$, it becomes an atom. Before acquisition the support consists only of the unresolved predictive point. The physical scales are
\[
 \psi_1(k)=a^2/k^2,\qquad
 \psi_2(k)=\max\{\max(b,c)^2/k^2,bc/k\}\quad(bc>0),
\]
with zero widths removed at the boundary. There is no single smooth acquired manifold at the crossing or a single ambient-dimension exponent. Theorem~\ref{thm:sensor} applies with constants uniform through all these changes. Smooth nonalgebraic bi-Lipschitz distortions of the two charts yield the same conclusion, with their uniform metric constants. Thus the extension is not a statement about a preassigned Hilbert-space rank cut.

\subsection{Recurrent expansion with intermittent refresh}\label{sec:recurrent}
We give a third realization to test the structural kernel theorem, rather than only its contractive and absorbing special mechanisms. Fix $l=1/4$, $0\le a\le1/2$, and $u=l+a$. Prepare $X_0$ uniformly on $[l,u]$; for $a=0$ this is a point. One initial paid call reveals $X_0$. At each subsequent paid tick, independently of the past, a hold occurs with probability $1-\theta$, $0\le\theta\le1$. On a hold, $X$ is unchanged and the report is the symbol $H$. Otherwise an active type is chosen independently: with probability $p=1/10$ it is the fold $E$, and with probability $1-p$ it is the refresh $C$. The raw transition/report kernels are
\begin{align}\label{eq:recurrent-kernel}
 E:&\quad X'=l+2\min\{X-l,u-X\},\quad Y=E,\\
 C:&\quad X'=cX+(1-c)(l+aU),\quad Y=(C,U),
           \quad c=1/4,\quad U\sim\operatorname{Unif}[0,1].\notag
\end{align}
These are normalized nonnegative Borel kernels; the fold is deterministic. There is no further exact reveal of $X'$. A terminal probe is Bernoulli with mean the current $X$. All ticks, including the initial reveal and holds, cost one raw call. The exploration has no choice of informative actions and is fixed. Its protocol clock is exogenous.

\begin{theorem}[Acquired geometry under recurrent expanding updates]\label{thm:recurrent}
For every $N\ge1$, $m\ge1$, $a\in[0,1/2]$ and $\theta\in[0,1]$, the raw experiment \eqref{eq:recurrent-kernel} has predictive quotient $[l,u]$ with its protocol interface, and
\begin{equation}\label{eq:recurrent-risk}
 R_{\cp}(N,m)-B_N\asymp R_{\on}(N,m)-B_N\asymp a^2m^{-2},
 \qquad B_N=\E[X_N(1-X_N)].
\end{equation}
Here $X_N$ denotes the state after the $N$ paid calls, including the initial reveal. The constants are absolute and uniform in all displayed parameters. At $a=0$ both exact excesses are zero. The online bound has the kernel-derived transport constant
\begin{equation}\label{eq:recurrent-gamma}
 \Gamma\le\frac{3040}{667},
\end{equation}
including initial rounding, independently of the hold probability.

If current inputs, physical endpoints and arithmetic have absolute errors controlled by $\eta$ per initial or active update, and hold operations are exact, the finite-bit implementation has
\[
 R_{\rm alg}(N,m)-B_N\le C\bigl(a/m+\eta\bigr)^2.
\]
The upper is independent of $N$ and $\theta$. The endpoint/error guarantee includes a projection back to the calibrated valid interval with its stated absolute accuracy. Its program and workspace costs are those of Appendix~\ref{sec:algorithms}; no stored exact real is used.
\end{theorem}
\begin{proof}
After the initial report the exact state is known. The reported type and, at refresh, the reported $U$, determine each subsequent exact update. All future laws are therefore functions of the current $x$ and the protocol interface. The terminal Bernoulli mean is $x$, so distinct states are predictively distinct. Continuation laws and these explicit updates are continuous on the interval within each report component. The report law before the terminal probe is independent of $x$, hence has zero report-kernel continuity constant. The terminal report has total variation $|x-x'|$. This establishes the quotient and report certificates directly from the raw kernel, without assuming a covariance geometry.

For $a>0$, let $f$ be the actual current density. Initially $f=1/a$. A hold preserves it. The tent fold has two affine inverse branches, each with Jacobian factor $1/2$, so its pushforward density is the average of the two branch densities and has sup norm at most $\|f\|_\infty$. For refresh, condition on the preceding state. The new state is uniform on an interval of length $(1-c)a$; its conditional density is at most $1/((1-c)a)=4/(3a)$. Mixing over the actual preceding law preserves this bound. Induction, including mixing over the reported type and holds, gives
\[
 0\le f_N(x)\le4/(3a),\qquad \int_l^u f_N(x)\,dx=1
\]
for every $N\ge1$. Thus the actual-mass certificate has $C_g=4/3$ on one chart of physical width $a$. It is a bound on the law generated by the experiment, not an imposed invariant uniform law. The task map is $x\mapsto x$ and $\psi(m)=a^2/m^2$.

Use a midpoint grid with $m$ representatives and physical Euclidean distance. Execute holds without arithmetic or rounding. Execute an active update from the previous representative and then round. The fold has Lipschitz constant two; refresh has constant $c=1/4$. Their rounding error is at most $a/(2m)$. Therefore the error recurrence has active gain $G=2$ with probability $1/10$ and $G=1/4$ with probability $9/10$, with forcing one after scaling the grid radius. The coefficients are independent marks, independent of the grid. In the one-state environmental kernel,
\[
 \E G=17/40,\qquad \E G^2=73/160<1.
\]
Theorem~\ref{thm:kernel} and \eqref{eq:scalar-resolvent} yield
\[
 H=\frac{1+17/40}{(1-17/40)(1-73/160)}=\frac{3040}{667}.
\]
The normalized initial rounding seed is at most one; the scalar fixed-point first and second moments are both at least one, so it is an admitted seed. Proposition~\ref{prop:suspension} proves that insertion of holds leaves this bound unchanged for $\theta>0$. At $\theta=0$ only the initial rounding persists, so the same upper holds directly. Apply Lemma~\ref{lem:initialization} to the one paid initial reveal and the $N-1$ subsequent calls. Theorem~\ref{thm:main} then supplies \eqref{eq:recurrent-risk}, and its lower uses the actual density just proved. No realization calculation has entered the proof of the kernel theorem.

On an active or initial update, interval projection and approximate arithmetic add at most $C\eta$ to the same error recurrence. On holds the representative is copied exactly, giving zero new error. The same scalar forcing bound, applied to radius $a/(2m)+C\eta$, proves the finite-bit assertion. The physical formulas in \eqref{eq:recurrent-kernel} use $a$ only multiplicatively; their absolute-error bounds never divide by it. At $a=0$ there is a single exact state, and the exact program can bypass all rounding. This also proves rank-uniformity.
\end{proof}

For every $\theta>0$ the active folds occur repeatedly, and blocks of consecutive folds have positive probability. Thus this example is not uniformly contractive along active paths, not an absorbing reveal after the initial preparation, and not a finite measurement ensemble. Its stability is a forced second-moment conclusion derived from the raw update probabilities. Uniformity in arbitrarily long waiting periods does not remove those periods from the acquisition or physical-time budget.
